% Part: computability
% Chapter: computability-theory
% Section: computable-sets

\documentclass[../../../include/open-logic-section]{subfiles}

\begin{document}

\olfileid[hi]{cmp}{thy}{cps}
\olsection{संगणनीय समुच्चय}

हम संगणनीयता की धारणा को संगणनीय फलनों से संगणनीय समुच्चयों तक विस्तृत
कर सकते हैं:

\begin{defn}
  $S$ को प्राकृतिक संख्याओं का समुच्चय मानें। तब $S$ \emph{संगणनीय}
  तभी और केवल तभी है जब उसका अभिलाक्षणिक फलन~$\Char{S}$ संगणनीय हो।
  दूसरे शब्दों में, $S$~तभी और केवल तभी संगणनीय है जब फलन
\[
\Char{S}(x) =
\begin{cases}
1 & \text{यदि $x \in S$} \\
0 & \text{अन्यथा}
\end{cases}
\]
संगणनीय हो। इसी प्रकार संबंध $R(x_0, \dots, x_{k-1})$ तभी और केवल तभी
संगणनीय है जब उसका अभिलाक्षणिक फलन संगणनीय हो।

संगणनीय समुच्चयों और संबंधों को \emph{निर्णेय} भी कहा जाता है।
\end{defn}

\begin{explain}
ध्यान दें कि अब हमारे पास संगणनीयता की कई धारणाएँ हैं: आंशिक फलनों,
फलनों और समुच्चयों के लिए। इन्हें परस्पर न मिलाएँ!{}
\iftag{TMs}{किसी आंशिक फलन को संगणित करने वाली ट्यूरिंग मशीन उन आगत
  मानों पर फलन का निर्गत लौटाती है जिन पर फलन परिभाषित है; किसी समुच्चय
  को संगणित करने वाली ट्यूरिंग मशीन यह तय करने के बाद कि आगत मान समुच्चय
  में है या नहीं, या तो $1$ अथवा~$0$ लौटाती है।}{}
\end{explain}

\end{document}
